\scanpage{53}
\diagram{\begin{tikzcd}[ampersand replacement=\&,column sep=small]
 I\arrow[d,"p"] \& I_\alpha\\ A \& \alpha
\end{tikzcd}}
$S\simeq$ ens des parties \inserted{$\deleted{\unsure{I}}A$} de cardinal $2$ de $I$, telles que $p|A$ soit injectif. (sections partielles sur parties de cardinal $2$)

\figure{\begin{tikzpicture}[scale=.44]
\foreach \x in {0,1,2}{\foreach \y in {1.5,2.5,3.5}{\fill (\x,\y) circle (.042);}\fill (\x,0) circle (.065);}
\draw[line width=1.6pt] (0,3.5)--(1,2.5);
\draw[->] (2.75,1.2)--(2.75,.34);
\end{tikzpicture}}
$X_\alpha=$ ens des $A$ tels que $p(A)=A-\{\alpha\}$

$A,B\in S$ liés ssi $A\ne B$ et
\[
\begin{aligned}
 &\text{ou bien }p(A)=p(B)\qquad A\cap B=\varnothing\\
 &\text{ou bien }p(A)\ne p(B),\quad A\cap B\ne\varnothing
\end{aligned}
\]

\textcircled{\scriptsize 12} La catégorie des graphes cubiques, munis d'une décomposition en trois \deleted{\unsure{bistriangles|bistriples|bitriangles}} \inserted{bitriangles} disjoints, est équivalente à \deleted{celle} \inserted{la cat. $(\mathrm{Ens})_{(3,3,3)}$} des ens \deleted{de type} munis d'une partition de type $(3,3,3)$. Le groupe des autom d'une telle structure est $\simeq\mathfrak S_3\mathbin{.}(\mathfrak S_3\times\mathfrak S_3\times\mathfrak S_3)$ qui est d'ordre $3^4\,2^4$ (c'est donc \inserted{presque} un sous-groupe de Sylow de $\mathrm{Aut}(\Delta)\deleted{\unsure{\ne2}}$ !)
\inserted{il faut se donner des \deleted{\unsure{orientations}} \inserted{\unsure{ordres|cycles|cinq} circ.} des $I_\alpha$ et de $A$ et exiger qu'ils soient respectés \ldots}

% slanted
\margin{\uline{NB} La \deleted{\unsure{réunion|réunition}} \inserted{partition} de $S$ en les $X_\alpha$ se reconstitue, à partir de ce ss-gpe de $\mathrm{Aut}(\Delta)$, comme l'ens des $3$ orbites sur $S$ de ce gpe: donc ce ss-gpe de Sylow est son propre normalisateur.}

La cat. des couples $(\Delta,X)$ d'un graphe cubique, et d'un bitriangle $X$ dans $\Delta$, est équivalente à \deleted{celle des} \inserted{la cat.} \deleted{couples} $(\mathrm{Ens})_3\times(\mathrm{Ens})_{3,3}$ des couples formés d'un ens de card. $3$ (savoir l'ens \deleted{\unsure{des élémts}} \inserted{can. de} \unsure{$P$|$P$ de} orbites associé à $X$ dans $\Delta$) et d'un ens à partition de type $(3,3)$ (savoir l'ens des triades de $X$, \deleted{\unsure{\}|)}} avec sa partition naturelle).
% archive sheet number: 52
